\hypertarget{uxb300uxc0c1-uxc804uxc81cuxc640-uxd45cuxae30}{%
\section{대상, 전제와 표기}\label{uxb300uxc0c1-uxc804uxc81cuxc640-uxd45cuxae30}}

같은 초기 COMMON과 하나의 admissible 고정 계수 선택 \(\alpha\)에 대해 두 전역 매개변수 \(\theta=(\lambda,b)\in[0,1]^2\)를 Full과 실제 두 half 경로에 공유한다. 미분할 매개변수는 \(\theta\in\{\lambda,b\}\) 하나이며 점은 도트로도 표기한다. 시간, 밀도, 원 상수, 수송 사상, COMMON의 미래 매개변수 도함수는 고정한다. 원 source family의 작은 열린 매개변수 확장 덕분에 닫힌 rectangle의 경계에서도 해당 1차 미분을 정의한다.

기체 좌표와 endpoint는

\[
g=(x,y,z,w)^T,\quad y_*=g_2,\quad d=g_2-g_F,\quad
v=g_{1,\theta},\quad u=g_{2,\theta},\quad s=g_2-g_1
\]

로 둔다. 이하 \(y_*\)는 두 번째 endpoint 벡터이고, \(y\)는 He II fraction이다. \(w\)의 단위는 eV/H이며 \(\lambda,b\)는 무차원이다. \(h=625000000\ \mathrm{s}\), Full 길이는 \(2h\)이다. 동일 endpoint 시각의 정확한 density ratio를 \(f=f_2=\operatorname{exact}(n_{\mathrm{He},2}^{\mathrm{stored}})/\operatorname{exact}(n_2)\)라 둔다. 물리상수 \(c,k_B\) 및 eV-to-erg leaf를 원 값으로 유지한다.

사용하는 가정은 다음과 같다.

\begin{enumerate}
\def\labelenumi{\arabic{enumi}.}
\tightlist
\item
  원 4차원 convex \(X\)의 내부 hull에 \(g_1,g_2,g_F\)와 필요한 선분·density pullback이 포함된다. source와 실제 implicit root family는 상속된 \(C^2\) family이다.
\item
  같은 고정 \(C,B,S,q\)가 \(X\)에서 \(\lvert I-CH_g(g;\lambda,b)\rvert\le B\), \(\lVert S^{-1}BS\rVert_\infty\le q<1\)을 만족한다. \(H=H_F\)는 Full 잔차다.
\item
  PHYS09/11/14의 원 1차 gas·photon carry와 PHYS16의 상관된 \(d,s\), 잔차 포함을 재사용한다. 원 signed gradient/Hessian은 부호 clipping을 하지 않는다.
\item
  현재 nonphoto source는 원 FT03의 CI·RR·DR graph이다. HI threshold \(\chi_H\)는 실제 source의 동일 captured leaf이다. 추가 source 성분은 별도 \(B\) 항으로 명시한다.
\end{enumerate}

이 가정과 원본 identity는 PHYS16의 \nolinkurl{contracts/PARAMETER_RESPONSE_CONTRACT.json}, PHYS15의 \nolinkurl{theory/ASSUMPTIONS_AND_BINDINGS_KO.md} 및 PHYS16의 \nolinkurl{evidence/intake/SOURCE_COOLING_AUDIT_KO.md}에 연결된다. 다음 유도에서 정칙성은 \(C^2\) 안에 유지된다.

원 source 포함의 적용 위치는 아래와 같다. 새로운 연결식이 원 \(X\) 밖의 source 값을 요구하지 않는다.

\begin{longtable}[]{@{}
  >{\raggedright\arraybackslash}p{(\columnwidth - 2\tabcolsep) * \real{0.5000}}
  >{\raggedright\arraybackslash}p{(\columnwidth - 2\tabcolsep) * \real{0.5000}}@{}}
\toprule\noalign{}
\begin{minipage}[b]{\linewidth}\raggedright
식의 평가점
\end{minipage} & \begin{minipage}[b]{\linewidth}\raggedright
도함수 입력과 포함 근거
\end{minipage} \\
\midrule\noalign{}
\endhead
\bottomrule\noalign{}
\endlastfoot
\(g_1,g_2\) 및 \(g_1+ts\) & 원 first-half/실제 second-half root hull의 convex hull \(\subset X\); Full nonphoto/q box Hessian과 centre-to-hull gradient 포함 \\
\(g_1\)의 온도 정렬점 \(g_*=\Psi(g_1)\)와 \(g_1+t(g_*-g_1)\) & PHYS15에서 exact density ratio와 내부 hull로 검토한 pullback; PHYS16 적용에서 같은 값을 재사용 \\
\(g_2-td,\ 0\le t\le1\) & PHYS16에서 포함한 실제 \(d=g_2-g_F\)와 원 \(g_2,g_F\)의 공통 convex \(X\) \\
P2의 \(H_{gg}\)와 \(H_{g\theta}\) & 원 Full \texttt{SOURCE\_BOX.json}의 잔차 gas Hessian 및 mixed gas/parameter Hessian; \(\lambda\) slot 4, b slot 5 \\
\end{longtable}

centre 자료를 이용해 값·gradient 포함을 좁힐 때도 Hessian의 권위 영역은 원 \(X\)이다. Whole-X Hessian을 새 centred Hessian 값으로 바꾸거나 추가 smoothness를 가정하지 않는다.

\hypertarget{uxc815uxb9ac-1-uxc791uxc740-endpoint-uxcc28uxc774uxb97c-uxba3cuxc800-uxc778uxc218-uxbd84uxd574uxd55c-uxbbf8uxbd84uxc2dd}{%
\section{정리 1 --- 작은 endpoint 차이를 먼저 인수 분해한 미분식}\label{uxc815uxb9ac-1-uxc791uxc740-endpoint-uxcc28uxc774uxb97c-uxba3cuxc800-uxc778uxc218-uxbd84uxd574uxd55c-uxbbf8uxbd84uxc2dd}}

\[
H(g)=g-g_0-2hB_2(g,\lambda)-VL_F(g)Q_F(b),\quad
H(g_F)=0,\quad r=H(g_2).
\]

\[
J_y=H_g(g_2),\qquad J_F=H_g(g_F),
\]

\[
\mathcal L(d)v_0=\int_0^1H_{gg}(g_2-td)[d,v_0]dt,\qquad
\mathcal G_\theta=\int_0^1H_{g\theta}(g_2-td)dt
\]

라 두면 정확히

\[
J_y-J_F=\mathcal L(d),\quad
H_\theta(g_2)-H_\theta(g_F)=\mathcal G_\theta d,
\]

\[
\boxed{J_Fd_\theta=f_\theta,
\qquad f_\theta=r_\theta-\mathcal L(d)u-\mathcal G_\theta d.}
\tag{1}
\]

\textbf{증명.} \(H(g_F;\theta)=0\)을 미분하면 \(J_F(g_{2,\theta}-d_\theta)+H_\theta(g_F)=0\)이다. 한편 \(r_\theta=J_yu+H_\theta(g_2)\)다. 두 식을 빼고 \(H_g\)와 \(H_\theta\)에 각각 선분 적분을 적용하면 (1)이 성립한다. \(H_{gg},H_{g\theta}\)의 연속성만 사용한다. \(\square\)

따라서 새 forcing \(f_\theta\)의 모든 feedback 항은 작은 \(d\)를 명시적으로 인수로 갖는다. 이와 동치인 endpoint 기준식은

\[
\boxed{J_yd_\theta=f_\theta+\mathcal L(d)d_\theta.}
\tag{2}
\]

원 Full 기체 1차 carry \(g_{F,\theta}\)를 사용할 때에는

\[
J_yd_\theta=r_\theta-\mathcal L(d)g_{F,\theta}-\mathcal G_\theta d
\tag{3}
\]

도 같은 \(d_\theta\)를 포함한다. 식 (3)은 두 경로의 원 1차 carry와 상관 잔차의 독립 연결로 사용할 수 있다.

\hypertarget{uxc2e4uxc81c-source-slotuxacfcuxc758-uxb300uxc751}{%
\subsection{실제 source slot과의 대응}\label{uxc2e4uxc81c-source-slotuxacfcuxc758-uxb300uxc751}}

현재 \(H\)의 명시적 매개변수 미분은

\[
H_\lambda=-2h q_2\nu_H,\qquad
H_{g\lambda}=-2h\nu_H\otimes Dq_2,
\]

\[
H_b=-2VL_FJ_\gamma,\qquad
H_{gb}=2\sum_j\frac{\ell_jJ_{\gamma j}}{D_{Fj}^2}\nu_j e_x^T,
\quad\ell_j=2ha_{2j}.
\tag{4}
\]

여기서 \(\nu_H=(1,0,0,-\chi_H)^T\), \(J_\gamma=M_\gamma e_{24}\), \(D_{Fj}=1+\ell_j(1-x)\)이다. source export의 \texttt{G\_g\_lambda{[}i{]}{[}j{]}}, \texttt{G\_g\_b{[}i{]}{[}j{]}}는 각각 residual{[}i{]}.hessian{[}j{]}{[}4{]}, residual{[}i{]}.hessian{[}j{]}{[}5{]}에 대응한다. \texttt{G\_gg}는 이미 Full \textbf{잔차} Hessian이다.

\hypertarget{uxc815uxb9ac-2-stage-movementuxc758-uxb3c4uxd568uxc218uxb97c-uxacf5uxc720uxd558uxb294-gas-uxc794uxcc28}{%
\section{정리 2 --- stage movement의 도함수를 공유하는 gas 잔차}\label{uxc815uxb9ac-2-stage-movementuxc758-uxb3c4uxd568uxc218uxb97c-uxacf5uxc720uxd558uxb294-gas-uxc794uxcc28}}

비광자 source를

\[
B_k(g,\lambda)=F_k(g)+\lambda q_k(g)\nu_H-2H_kw e_w
\]

로 쓴다. \(F_k\)는 원 nonphoto export이다. PHYS15/16에서

\[
s=hB_2(g_2)+VP_2,\qquad
r_g=h\{B_1(g_1)-B_2(g_2)\}
\]

를 얻었다. 같은 매개변수로 미분하면

\[
\boxed{s_\theta=h\{B_{2,g}(g_2)u+\delta_{\theta\lambda}q_2(g_2)\nu_H\}+VP_{2,\theta}.}
\tag{5}
\]

이 \(s_\theta\)는 정확히 \(u-v\)와 같으므로 보존된 두 sensitivity 포함과 교차할 수 있다. \(g_t=g_1+ts\), \(\delta B_t(g)=B_1(g)-B_2(g)\), \(\delta q_t=q_1-q_2\)로 두면

\[
\boxed{\begin{aligned}
r_{g,\theta}={}&h\{D\delta B_t(g_1)[v]+\delta_{\theta\lambda}\delta q_t(g_1)\nu_H\}\\
&-hB_{2,g}(g_2)s_\theta
-h\int_0^1B_{2,gg}(g_t)[s,v]dt
-h\delta_{\theta\lambda}\int_0^1Dq_2(g_t)[s]dt\,\nu_H.
\end{aligned}}
\tag{6}
\]

\textbf{증명.} 정확한 직접 미분은

\[
r_{g,\theta}=h\{B_{1,g}(g_1)v-B_{2,g}(g_2)u
+\delta_{\theta\lambda}[q_1(g_1)-q_2(g_2)]\nu_H\}.
\]

\(u=v+s_\theta\)를 넣는다. \(B_{1,g}(g_1)-B_{2,g}(g_2)\)는 \(D\delta B_t(g_1)-\int_0^1B_{2,gg}(g_t)[s,\cdot]\,dt\)이고, \(q_1(g_1)-q_2(g_2)=\delta q_t(g_1)-\int_0^1Dq_2(g_t)[s]\,dt\)다. 대입하면 (6)이 따른다. \(\square\)

식 (6)은 평균 source Jacobian 자체를 매개변수 미분한 표현을 endpoint \(B_{2,g}(g_2)s_\theta\)와 작은 \(s\)의 Hessian contraction으로 재구성한다. 원 \(H_1=H_2\)이므로 \(\delta B_t\)의 팽창항은 정확히 소거되고, \(B_{2,g}\)에는 \(-2H_2e_we_w^T\)를 유지한다.

\hypertarget{uxc815uxb9ac-3-cuxb2-nonphoto-jetuxc5d0uxc11c-he-thermal-derivativeuxb97c-uxc9c1uxc811-uxd68cuxc218}{%
\section{정리 3 --- \texorpdfstring{\(C^2\)}{C2} nonphoto Jet에서 He thermal derivative를 직접 회수}\label{uxc815uxb9ac-3-cuxb2-nonphoto-jetuxc5d0uxc11c-he-thermal-derivativeuxb97c-uxc9c1uxc811-uxd68cuxc218}}

이 정리는 density pullback의 작은 He 보정에 필요한 도함수를 닫는다. \(f,n\)은 endpoint 2의 고정 값이고

\[
e=x+f(y+2z),\quad \Pi=1+f+e,\quad T=C_Tw/\Pi,\quad
\ell_e=(1,f,2f,0)
\]

이다. 다음 정규화된 화학·열 조합과 온도를 유지하는 \(x\) 방향을 정의한다.

\[
U(g)=\frac{F_w(g)+\chi_HF_x(g)}{ne},\qquad
\xi(g)=e_x+\frac{w}{\Pi}e_w.
\tag{7}
\]

원 source에서

\[
F_x=ne\{(1-x)\beta_H-x\alpha_H\},
\]

\[
F_w=ne\{-\chi_H(1-x)\beta_H-x\kappa_H
+f k_{w\mathrm{He}}(y,z,T)\}.
\]

따라서

\[
U(g)=-x\{\kappa_H(T)+\chi_H\alpha_H(T)\}
+f k_{w\mathrm{He}}(y,z,T).
\tag{8}
\]

\(DT[\xi]=0\)이고 \(Dy[\xi]=Dz[\xi]=0\)이므로 식 (8)의 고정 온도 \(x\) 미분을 취하면

\[
DU[\xi]=-[\kappa_H(T)+\chi_H\alpha_H(T)].
\]

그러므로

\[
\boxed{k_{w\mathrm{He}}(g)=\frac{U(g)-xDU(g)[\xi(g)]}{f}.}
\tag{9}
\]

임의의 기체 방향 \(v_0\)에 대해 이를 한 번 미분하면

\[
\boxed{\begin{aligned}
Dk_{w\mathrm{He}}(g)[v_0]
=\frac{1}{f}\bigg\{
DU[v_0]-v_{0x}DU[\xi]
-x\big[D^2U[v_0,\xi]+DU[D\xi[v_0]]\big]\bigg\},\\
D\xi[v_0]=\left[\frac{v_{0w}}{\Pi}
-\frac{w\ell_ev_0}{\Pi^2}\right]e_w.
\end{aligned}}
\tag{10}
\]

\textbf{증명.} (8)은 원 hydrogen CI 항을 동일 \(\chi_H\)로 소거한 source 항등식이다. \(\xi\)를 따라 온도와 He fractions가 고정되므로 첫 방향 미분에서 수소 항의 계수만 남는다. (9)에 ordinary product·chain rule을 적용하고 \(\xi\)의 도함수를 계산하면 (10)을 얻는다. (9)는 현재 \(g\)에서의 국소 값·도함수 항등식으로, \(x=0\)까지 source를 평가할 필요가 없다. \(\square\)

\hypertarget{uxc800uxc7a5uxb41c-source-jetuxc5d0-uxbc14uxb85c-uxc801uxc6a9uxd558uxb294-uxacc4uxc0b0uxc2dd}{%
\subsection{저장된 source Jet에 바로 적용하는 계산식}\label{uxc800uxc7a5uxb41c-source-jetuxc5d0-uxbc14uxb85c-uxc801uxc6a9uxd558uxb294-uxacc4uxc0b0uxc2dd}}

\(L=F_w+\chi_HF_x\)라 두면 \(e\)의 Hessian은 0이므로

\[
DU[v_0]=\frac1n\left(\frac{DL[v_0]}e-
\frac{L\ell_ev_0}{e^2}\right),
\]

\[
\boxed{D^2U[v_0,z_0]=\frac1n\left\{
\frac{D^2L[v_0,z_0]}e
-\frac{DL[v_0]\ell_ez_0+DL[z_0]\ell_ev_0}{e^2}
+\frac{2L(\ell_ev_0)(\ell_ez_0)}{e^3}
\right\}.}
\tag{11}
\]

식 (9)--(11)의 입력은 Full \texttt{nonphoto{[}0{]}}, \texttt{nonphoto{[}3{]}}의 value, gas gradient, signed gas Hessian뿐이다. \(e,f,\Pi\)가 양수인 원 hull에서 모두 정칙하다. PHYS16의 cooling-cone 값 포함과 식 (9)의 값 포함은 같은 함수에 대한 두 포함으로 교차할 수 있다. 도함수 포함은 식 (10)--(11)에서 직접 얻는다.

\hypertarget{uxc815uxb9ac-4-exact-density-pullbackuxc758-uxc0c1uxad00uxb41c-1uxcc28-closure}{%
\section{정리 4 --- exact density pullback의 상관된 1차 closure}\label{uxc815uxb9ac-4-exact-density-pullbackuxc758-uxc0c1uxad00uxb41c-1uxcc28-closure}}

\[
\sigma=y+2z,\quad e_k=x+f_k\sigma,\quad
\Pi_k=1+f_k+e_k,\quad \Delta f=f_1-f_2,\quad r_n=n_1/n_2,
\]

\[
r_e=r_ne_1/e_2,
\quad\delta w=-w\Delta f(1+\sigma)/\Pi_1,
\quad g_*=g+\delta w e_w.
\]

그러면 \(T_2(g_*)=T_1(g)\)이고

\[
F_1(g)=r_eF_2(g_*)+n_1e_1\Delta f\,k_{w\mathrm{He}}(g_*)e_w,
\quad q_1(g)=r_nq_2(g_*).
\tag{12}
\]

작은 차이를 보존하는 값은

\[
r_e-1=(r_n-1)+r_n\Delta f\,\sigma/e_2.
\]

방향 \(v_0\)를 따라 \(\dot\sigma=v_{0y}+2v_{0z}\), \(\dot e_k=\ell_{e,k}v_0\)라 두면

\[
\boxed{Dr_e[v_0]=r_n\Delta f
\frac{e_2\dot\sigma-\sigma\dot e_2}{e_2^2},}
\tag{13}
\]

\[
\boxed{D\delta w[v_0]=-
\Delta f\left[
\frac{(1+\sigma)v_{0w}+w\dot\sigma}{\Pi_1}
-\frac{w(1+\sigma)\dot e_1}{\Pi_1^2}
\right].}
\tag{14}
\]

\(F=F_2\), \(q=q_2\), \(\delta F=F_1-F_2\), \(\delta q=q_1-q_2\), \(v_*=v_0+D\delta w[v_0]e_w\)로 줄여 쓰자. 정확히

\[
\boxed{\begin{aligned}
D\delta F(g)[v_0]={}&(r_e-1)DF(g)[v_0]+Dr_e[v_0]F(g_*)\\
&+r_e\delta w\int_0^1D^2F(g+t\delta w e_w)[e_w,v_0]dt\\
&+r_e(\partial_wF)(g_*)D\delta w[v_0]\\
&+n_1\Delta f\left\{
\dot e_1 k_{w\mathrm{He}}(g_*)
+e_1Dk_{w\mathrm{He}}(g_*)[v_*]
\right\}e_w,
\end{aligned}}
\tag{15}
\]

\[
\boxed{\begin{aligned}
D\delta q(g)[v_0]={}&(r_n-1)Dq(g)[v_0]\\
&+r_n\delta w\int_0^1D^2q(g+t\delta w e_w)[e_w,v_0]dt\\
&+r_n(\partial_wq)(g_*)D\delta w[v_0].
\end{aligned}}
\tag{16}
\]

여기서 \(\partial_wF\)와 \(\partial_wq\)는 \(w\) 좌표의 도함수다. 구현에서는 각각 \texttt{DF{[}:,3{]}}, \texttt{Dq{[}3{]}}로 명시한다.

\textbf{증명.} (12)를 \(v_0\)로 미분하고 \(Dg_*v_0=v_*\)를 사용한다. \(r_eDF(g_*)[v_0]-DF(g)[v_0]\)를 \((r_e-1)DF(g)[v_0]+r_e(DF(g_*)-DF(g))[v_0]\)로 나눈 뒤 마지막 차이에 Hessian 선분 적분을 적용한다. HH 식도 \(r_n\)이 고정인 같은 계산으로 따른다. He correction의 도함수는 (10)을 직접 대입한다. \(\square\)

식 (6)의 source-time 항은

\[
D\delta B_t(g)[v_0]=D\delta F(g)[v_0]+\lambda D\delta q(g)[v_0]\nu_H
\]

로 닫힌다. 현재 exact \(\Delta f\)는 약 \(-6.94525\times10^{-18}\)이고 이 작은 factor가 모든 density-normalization derivative 항에 유지된다. \(\Delta f=0\)이면 \(g_*=g\), \(D\delta w=Dr_e=0\)이므로 (15)--(16)은 정확히 \((r_n-1)DF[v_0]\), \((r_n-1)Dq[v_0]\)로 환원된다.

\hypertarget{uxc815uxb9ac-5-uxc6d0-photon-uxc77cuxacf1-sectoruxc758-1uxcc28-uxbbf8uxbd84}{%
\section{정리 5 --- 원 photon 일곱 sector의 1차 미분}\label{uxc815uxb9ac-5-uxc6d0-photon-uxc77cuxacf1-sectoruxc758-1uxcc28-uxbbf8uxbd84}}

\(R,R_F\)와 \(Z_0\)는 고정이고 \(J_\gamma=M_\gamma e_{24}\)이다. \(g_2\)에서

\[
\tau_j=ha_{2j}(1-x_2),\quad A=(I+\tau)^{-1},\quad
L=I-A,\quad A_F=(I+2\tau)^{-1},\quad L_F=I-A_F,
\]

\[
C_1=\tau A^2A_F,\quad C_2=\tau^2A^2A_F,\quad
Q_1=RZ_0+bJ_\gamma,\quad\Delta R=R^2-R_F.
\]

대각 연산자끼리의 곱은 성분별이다. 기체 전처리 \(C\)와 photon \(C_1,C_2\)는 구분한다. \(\delta_b=\delta_{\theta b}\)라 두면

\[
\dot Q_1=\delta_bJ_\gamma,\quad
\dot\tau_j=-ha_{2j}u_x,\quad
\dot A=-A^2\dot\tau,\quad \dot L=-\dot A,
\]

\[
\dot A_F=-2A_F^2\dot\tau,\quad \dot L_F=-\dot A_F,
\]

\[
\dot C_1=A^2A_F(1-2\tau A-2\tau A_F)\dot\tau,
\]

\[
\dot C_2=2\tau A^2A_F(1-\tau A-\tau A_F)\dot\tau.
\tag{17}
\]

첫 opacity 차이는 작은 scalar \(\zeta\)를 먼저 유지한다.

\[
\zeta_j=h\{c\sigma_j(n_2-n_1)(1-x_2)-a_{1j}s_x\},
\quad\Delta A_1=\zeta A_1A,
\]

\[
\dot\zeta_j=-h\{c\sigma_j(n_2-n_1)u_x+a_{1j}s_{\theta,x}\},
\quad \dot A_{1j}=ha_{1j}A_{1j}^2v_x,
\]

\[
\boxed{\dot{\Delta A}_1=\dot\zeta A_1A
+\zeta(\dot A_1A+A_1\dot A).}
\tag{18}
\]

\(K=I-LR\)라 두면 PHYS15의 각 \(p\) sector와 그 도함수는 다음과 같다.

\[
p_{op}=-K\Delta A_1Q_1,
\quad
\boxed{\dot p_{op}=\dot LR\Delta A_1Q_1
-K\dot{\Delta A}_1Q_1-K\Delta A_1\dot Q_1,}
\tag{19a}
\]

\[
p_{com}=L[R,A]Q_1,
\quad
\boxed{\dot p_{com}=\dot L[R,A]Q_1
+L[R,\dot A]Q_1+L[R,A]\dot Q_1,}
\tag{19b}
\]

\[
p_{loc}=L(I-R)RZ_0,
\quad \boxed{\dot p_{loc}=\dot L(I-R)RZ_0,}
\tag{19c}
\]

\[
p_{stock}=C_2R^2Z_0,
\quad \boxed{\dot p_{stock}=\dot C_2R^2Z_0,}
\tag{19d}
\]

\[
p_{split}=L_F\Delta RZ_0,
\quad \boxed{\dot p_{split}=\dot L_F\Delta RZ_0,}
\tag{19e}
\]

\[
p_{bR}=bLA(R-I)J_\gamma,
\quad
\boxed{\dot p_{bR}=\delta_bLA(R-I)J_\gamma
+b(\dot LA+L\dot A)(R-I)J_\gamma,}
\tag{19f}
\]

\[
p_{bBE}=-bC_1J_\gamma,
\quad
\boxed{\dot p_{bBE}=-\delta_bC_1J_\gamma-b\dot C_1J_\gamma.}
\tag{19g}
\]

\textbf{증명.} (17)은 각 대각 유리함수의 미분이다. (18)은 PHYS16의 동일 \(s\)에 의한 opacity 차이 식의 미분이고, (19)는 원 일곱 항에 product rule을 적용한 결과다. \(R\)과 \(A\)의 순서는 바꾸지 않는다. \(\square\)

식 (19f,g)는 \(b\)의 직접 항과 \(b\)에 따른 상태·opacity feedback을 함께 담는다. 전체 \(b\) 응답에서는 \(g_{1,b}\)와 \(g_{2,b}\)가 다른 모든 sector에도 들어간다.

\hypertarget{uxc2e4uxc81c-stage-derivativeuxc640uxc758-uxc815uxd655uxd55c-uxd569-uxac80uxc0ac}{%
\subsection{실제 stage derivative와의 정확한 합 검사}\label{uxc2e4uxc81c-stage-derivativeuxc640uxc758-uxc815uxd655uxd55c-uxd569-uxac80uxc0ac}}

\[
Q_2=RA_1Q_1+bJ_\gamma,\quad
\dot Q_2=R(\dot A_1Q_1+A_1\dot Q_1)+\delta_bJ_\gamma,
\]

\[
\dot P_1=-\dot A_1Q_1+(I-A_1)\dot Q_1,
\quad \dot P_2=\dot LQ_2+L\dot Q_2,
\]

\[
Q_F=R_FZ_0+2bJ_\gamma,\quad
\dot{\widehat P}_F=\dot L_FQ_F+2\delta_bL_FJ_\gamma.
\]

따라서

\[
\boxed{\sum_s\dot p_s=\dot P_1+\dot P_2-\dot{\widehat P}_F,
\qquad r_\theta=r_{g,\theta}+V\sum_s\dot p_s.}
\tag{20}
\]

33개 원 group과 number/energy guard 두 개에 그대로 적용한다. Guard에서 \(A=1,L=0\) 및 그 도함수 0을 유지한다.

\hypertarget{uxbbf8uxbd84uxb41c-photon-ledgeruxc640-uxc2e4uxc81c-endpoint-readout}{%
\subsection{미분된 photon ledger와 실제 endpoint readout}\label{uxbbf8uxbd84uxb41c-photon-ledgeruxc640-uxc2e4uxc81c-endpoint-readout}}

원 photon number·energy row를 각각 \(n,u_\gamma\)라 하고 \(\Delta\widehat Z=Z_2-A_F(g_2)Q_F\)로 두면

\[
\mathbf1^T\dot p=-\partial_\theta(n\Delta\widehat Z),
\]

\[
E^T\dot p=(1-\rho)\{E^T\dot P_1-\delta_bM_\gamma E_{24}\}
-\partial_\theta(u_\gamma\Delta\widehat Z).
\tag{21}
\]

원 stored heat \(\epsilon_j=E_j-\chi_H\)가 정확하므로 heat ledger는 위 두 식의 같은 선형 결합이다. 실제 Full endpoint에서는

\[
\Delta Z_j=\Delta\widehat Z_j+\Gamma_jd_x,
\quad \Gamma_j=\frac{\ell_jQ_{Fj}}{D_{yj}D_{Fj}},
\quad D_{Fj}=D_{yj}+\ell_jd_x,
\]

\[
\dot\Gamma_j=\frac{\ell_j\dot Q_{Fj}}{D_{yj}D_{Fj}}
+\Gamma_j\ell_j\left(\frac{u_x}{D_{yj}}+
\frac{u_x-d_{\theta,x}}{D_{Fj}}\right),
\]

\[
\boxed{\Delta Z_{j,\theta}=\Delta\widehat Z_{j,\theta}
+\Gamma_jd_{\theta,x}+\dot\Gamma_jd_x.}
\tag{22}
\]

Guard에는 이 opacity 보정이 0이다. 모든 식이 동일 \(d,d_\theta,u\)와 \(Q_{F,\theta}=2\delta_bJ_\gamma\)를 공유한다.

\hypertarget{uxc815uxb9ac-6-uxc6d0-cbqsuxc5d0-uxc758uxd55c-sensitivity-uxd3ecuxd568uxacfc-finite-neumann}{%
\section{정리 6 --- 원 C,B,q,S에 의한 sensitivity 포함과 finite Neumann}\label{uxc815uxb9ac-6-uxc6d0-cbqsuxc5d0-uxc758uxd55c-sensitivity-uxd3ecuxd568uxacfc-finite-neumann}}

식 (1)의 \(f_\theta\)를 구성한 뒤 원 전처리를 적용한다. \(E_F=I-CJ_F\)에 대해 \(\lvert E_F\rvert\le B\)이고

\[
d_\theta=Cf_\theta+E_Fd_\theta.
\]

비음수 \(m_\theta\)가

\[
\boxed{(I-B)m_\theta=\operatorname{mag}([Cf_\theta])}
\tag{23}
\]

를 만족하면

\[
\boxed{d_\theta\in[Cf_\theta]+[-Bm_\theta,Bm_\theta],
\quad \|d_\theta\|_S\le a_\theta:=\frac{\|Cf_\theta\|_S}{1-q}.}
\tag{24}
\]

\textbf{증명.} \(\lVert S^{-1}BS\rVert_\infty<1\)에서 \(I-B\)는 가역이고 역행렬은 비음수 Neumann 급수다. \(\lvert d_\theta\rvert\le\lvert Cf_\theta\rvert+B\lvert d_\theta\rvert\) 를 그 역행렬로 전달하면 (23)--(24)가 성립한다. \(\square\)

기존 \(d_\theta\) 포함 또는 \(g_{2,\theta}-g_{F,\theta}\)와 교차한 뒤

\[
D_\theta\leftarrow D_\theta\cap([Cf_\theta]+[E_F]D_\theta)
\]

로 signed 포함을 조일 수 있다. 이는 원 true endpoint 식에 대한 유한 포함 연산이다.

\hypertarget{uxb450-uxbc88uxc9f8-endpointuxb97c-uxae30uxc900uxc73cuxb85c-uxd55c-uxc624uxcc28uxac00-uxc788uxb294-uxc720uxd55c-uxbcf4uxc815}{%
\subsection{두 번째 endpoint를 기준으로 한 오차가 있는 유한 보정}\label{uxb450-uxbc88uxc9f8-endpointuxb97c-uxae30uxc900uxc73cuxb85c-uxd55c-uxc624uxcc28uxac00-uxc788uxb294-uxc720uxd55c-uxbcf4uxc815}}

\[
M_y=I-CJ_y,\qquad v_{m,\theta}=\sum_{k=0}^mM_y^kCf_\theta.
\]

원 PHYS16에서 C-weighted Hessian bound \(L_H\)와 \(a\ge\lVert d\rVert_S\)를 얻었다. 식 (1)의 \(\mathcal L\)에는 적분 weight가 1이므로

\[
\|C\mathcal L(d)\|_S\le L_Ha.
\]

식 (2)와 유한 geometric identity에서

\[
\boxed{d_\theta=v_{m,\theta}
+\sum_{k=0}^mM_y^kC\mathcal L(d)d_\theta
+M_y^{m+1}d_\theta,}
\tag{25}
\]

\[
\boxed{\|d_\theta-v_{m,\theta}\|_S
\le a_\theta\left[q^{m+1}+L_Ha\sum_{k=0}^mq^k\right].}
\tag{26}
\]

이 식의 Hessian 항에는 1/2가 없다. PHYS16의 primal Taylor 나머지는 \((1-t)\)를 적분하여 1/2가 생겼고, 여기의 Jacobian 차이는 weight 1을 적분하기 때문이다. \(J_F\) 자체를 기준으로 \(M_F=I-CJ_F\)와 유한 다항식을 사용하면 (26)의 \(L_Ha\) 항이 소거되어 오차는 \(q^{m+1}a_\theta\)이다.

원 PHYS16의 \(a\approx4.28320\times10^{-6}\), \(L_H\approx8.27803\times10^{-5}\)를 재사용하면 \(L_Ha\approx3.55\times10^{-10}\)이다. 실제 PHYS17 \(a_\theta\)와 \(m\)에 따른 오차는 적용 결과에서 별도로 평가한다.

\hypertarget{uxc815uxb9ac-7-uxc720uxd55c-uxc628uxb3c4-uxbbf8uxbd84uxacfc-uxbaa8uxb4e0-forcinguxc5d0-uxacf5uxd1b5uxc778-adjoint}{%
\section{정리 7 --- 유한 온도 미분과 모든 forcing에 공통인 adjoint}\label{uxc815uxb9ac-7-uxc720uxd55c-uxc628uxb3c4-uxbbf8uxbd84uxacfc-uxbaa8uxb4e0-forcinguxc5d0-uxacf5uxd1b5uxc778-adjoint}}

동일 endpoint \(f_2\)에서

\[
d_e=\ell_ed,\quad \Pi_F=\Pi_2-d_e>0,\quad
T_2=C_Tw_2/\Pi_2,\quad
\Delta T=(C_Td_w-T_2d_e)/\Pi_F.
\]

미분하면

\[
\boxed{\Delta T_\theta=
\frac{C_Td_{w,\theta}-T_2d_{e,\theta}
-T_{2,\theta}d_e-\Delta T(\ell_eu-d_{e,\theta})}{\Pi_F}.}
\tag{27}
\]

여기서

\[
T_{2,\theta}=\frac{C_Tu_w-T_2\ell_eu}{\Pi_2}.
\]

\(T_F=T_2-\Delta T\),

\[
n_F=C_Te_w^T-T_F\ell_e,\qquad
k_\theta=T_{2,\theta}d_e+\Delta T\ell_eu
\]

로 묶으면

\[
\boxed{\Delta T_\theta=(n_Fd_\theta-k_\theta)/\Pi_F.}
\tag{28}
\]

\textbf{증명.} 원 유한 온도 분수의 quotient rule을 적용하면 (27)이다. \(d_{e,\theta}=\ell_ed_\theta\)이며 \(T_2-\Delta T=T_F\)를 사용해 \(d_\theta\)에 곱해진 두 항을 모으면 (28)을 얻는다. \(\square\)

같은 true Full Jacobian의 adjoint \(z_F\)를

\[
J_F^Tz_F=n_F^T
\]

로 정의하면

\[
\boxed{\Delta T_\theta=(z_F^Tf_\theta-k_\theta)/\Pi_F.}
\tag{29}
\]

\(f_\theta\)를 직접 forcing·gas 시간 변화·movement·일곱 photon derivative·endpoint feedback으로 분해했을 때, 모든 항에 \textbf{같은 \(J_F,n_F,\Pi_F\)}를 사용한다. 마지막 \(k_\theta\)는 관측량 자체의 상태 의존성에서 오는 공통 보정이다.

finite adjoint row를

\[
z_{m,\theta}^T=n_F\sum_{k=0}^mM_y^kC
\]

로 두고 \(\beta_F=\lVert n_FS\rVert_1\), \(\Pi_F\ge\pi>0\)라 하면

\[
\boxed{\left|\Delta T_\theta-
\frac{z_{m,\theta}^Tf_\theta-k_\theta}{\Pi_F}\right|
\le\frac{\beta_Fa_\theta}{\pi}
\left[q^{m+1}+L_Ha\sum_{k=0}^mq^k\right].}
\tag{30}
\]

각 forcing \(f_{\theta,s}\)의 endpoint response \(d_{\theta,s}=J_F^{-1}f_{\theta,s}\)에는 \(a_{\theta,s}=\lVert Cf_{\theta,s}\rVert_S/(1-q)\)를 써서 같은 식을 적용한다. 전체 합을 이루는 정확한 정의는 \(\Delta T_\theta=\sum_s(n_Fd_{\theta,s})/\Pi_F-k_\theta/\Pi_F\)다.

원 \(C\) 자체가 들어가는 row 차이도 명시할 수 있다. \(c_C=\lVert S^{-1}CS\rVert_\infty\), \(\kappa=(1-q)^{-1}\)라 두면

\[
\boxed{\|(z_F^T-z_{m,\theta}^T)S\|_1
\le\beta_F\kappa c_C
\left[q^{m+1}+L_Ha\sum_{k=0}^mq^k\right].}
\tag{31}
\]

\textbf{증명.} (25)에 \(n_F\)를 적용하고 dual scaled norm을 사용하면 (30)이 따른다. 행렬 식 \(J_F^{-1}-\sum_{k=0}^{m}M_y^kC=[M_y^{m+1}+\sum_{k=0}^{m}M_y^kC\mathcal L]J_F^{-1}\)에 \(\lVert S^{-1}J_F^{-1}S\rVert_\infty\le\kappa c_C\)를 적용하면 (31)을 얻는다. \(\square\)

\hypertarget{uxadf9uxd55c-uxcc28uxc6d0uxacfc-uxacfcuxd559uxc801-uxd310uxbcc4}{%
\section{극한, 차원과 과학적 판별}\label{uxadf9uxd55c-uxcc28uxc6d0uxacfc-uxacfcuxd559uxc801-uxd310uxbcc4}}

\(\theta\)가 고정 coefficient와 COMMON을 바꾸지 않는다는 전제가 모든 미분의 기준이다. 수송·density·시간에 대한 추가 매개변수를 도입할 경우 그 직접 도함수를 식 (12),(17),(19)에 추가한다.

\begin{itemize}
\tightlist
\item
  \textbf{같은 source 시각 극한:} \(n_1=n_2\)와 \(f_1=f_2\)이면 (15),(16)의 source-time derivative가 0이다.
\item
  \textbf{투명 광자 극한:} 모든 \(a_{kj}=0\)이면 \(A=A_1=A_F=I\), \(L=L_F=C_1=C_2=0\)이고 absorption sector와 그 매개변수 미분은 0이다.
\item
  \textbf{\(b=0\):} birth sector 값은 0이지만 \(b\) 미분의 \(\delta_b\) 항은 남는다. 이는 총 \(b\) 응답을 구성하는 직접 injection 방향이다.
\item
  \textbf{\(\lambda\) 미분:} \(\delta_{\theta b}=0\)이어도 원 source가 만든 \(g_{1,\lambda},g_{2,\lambda}\)를 통해 photon absorption과 모든 opacity feedback에 \(\lambda\) 응답이 들어간다.
\item
  \textbf{endpoint 일치:} \(d=0\)이면 (1)의 두 correction은 0이다. parameter family 전체에서 두 경로가 같고 \(r_\theta=0\)일 때 \(d_\theta=0\)이다.
\item
  \textbf{단위:} \(r_\theta\),\(f_\theta\),\(d_\theta\)의 앞 세 성분은 fraction, 마지막은 eV/H이고 모두 단위 무차원 parameter당 값이다. \(B_g\)가 \(\mathrm{s}^{-1}\)을 포함하므로 \(hB_gs_\theta\)는 기체 단위다. 식 (15)의 \(n\Delta fDk_{w\mathrm{He}}\)는 eV/H/s이고, (30)의 bound는 K/parameter다.
\end{itemize}

이번 정리의 실제 판별 대상은 전체 \(\partial_\lambda\Delta T\)와 \(\partial_b\Delta T\)의 포함·부호 여유다. 판단에 필요한 원 source·root data가 이미 있어, (5),(6),(10),(15)--(20),(1),(24),(27)을 순서대로 적용할 수 있다. PHYS16의 sector별 온도 기여가 이번에는 source forcing과 경로 feedback을 모두 가진 실제 매개변수 응답으로 연결된다.

\hypertarget{uxc2e4uxc81c-1uxcc28-uxd3ecuxd568uxc5d0uxc11c-uxc5bbuxb294-uxbaa8uxc11cuxb9acuxc720uxd55c-uxbcc0uxd654uxb4f1uxacb0uxd568-uxace1uxc120}{%
\subsection{실제 1차 포함에서 얻는 모서리·유한 변화·등결함 곡선}\label{uxc2e4uxc81c-1uxcc28-uxd3ecuxd568uxc5d0uxc11c-uxc5bbuxb294-uxbaa8uxc11cuxb9acuxc720uxd55c-uxbcc0uxd654uxb4f1uxacb0uxd568-uxace1uxc120}}

\begin{flushleft}
\nolinkurl{checks/application/run02/results/PARAMETER_RESPONSE_RESULTS.json}의 exact rational endpoint를 사용한다. 원 파일은 626092 bytes, SHA256 \texttt{fd4fc0fe6bcf0354\allowbreak{}96270bb92f0c56c4\allowbreak{}4ca43d16694cb402\allowbreak{}96087569a305e7c0}이다. 이하 corollary는 이 원 포함의 직접 결과이며, 실제 채택 상태는 독립 review 기록에 연결된다. 안전한 외향 표기는
\end{flushleft}

\[
\ell_-\le\partial_\lambda\Delta T\le\ell_+<0,
\quad [\ell_-,\ell_+]=[-1.009599125635,-0.8137506081166]\,10^{-8}\ {\rm K},
\]

\[
0<\beta_-\le\partial_b\Delta T\le\beta_+,
\quad[\beta_-,\beta_+]=[6.776348385589,7.185335549648]\,10^{-5}\ {\rm K}
\tag{32}
\]

이다. 하나의 \(\alpha\)를 고정하면 rectangle에서

\[
\boxed{\Delta T(1,0)\le\Delta T(\lambda,b)\le\Delta T(0,1).}
\tag{33}
\]

또 \(\Delta\lambda,\Delta b\ge0\)인 두 점이 원 rectangle에 있으면

\[
\boxed{\ell_-\Delta\lambda+\beta_-\Delta b
\le\Delta T(\lambda+\Delta\lambda,b+\Delta b)-\Delta T(\lambda,b)
\le\ell_+\Delta\lambda+\beta_+\Delta b.}
\tag{34}
\]

특히 각 parameter를 0에서 1로 바꾼 유한 변화는 (32)의 해당 도함수 구간에 포함된다. \(\lambda\)를 증가시켜 얻는 변화는 약 −10.10에서 −8.14 nK이고, \(b\)를 증가시켜 얻는 변화는 약 +67.76에서 +71.86 \(\mu\mathrm K\)이다. 고정 \(\alpha\)에서 두 매개변수가 만드는 총 범위는 exact endpoint 적분으로

\[
\boxed{\Delta T(0,1)-\Delta T(1,0)
\in[6.777162136197765,7.186345148772871]\,10^{-5}\ {\rm K}.}
\tag{35}
\]

원 내부 점 \((\lambda_0,b_0)\)을 지나는 \(\Delta T(\lambda,b)=\Delta T(\lambda_0,b_0)\)는 \(\partial_b\Delta T>0\)에 의해 국소적으로 유일한 \(C^1\) 함수 \(b=b_*(\lambda)\)를 정의하고

\[
\boxed{\frac{db_*}{d\lambda}=-\frac{\partial_\lambda\Delta T}{\partial_b\Delta T}
\in[1.132515805968976,1.489886688502190]\,10^{-4}>0.}
\tag{36}
\]

예를 들어 \((1-\lambda_0)\cdot1.489886688502190\times10^{-4}<1-b_0\)이면 \(\lambda_0\)에서 \(\lambda=1\)까지 그 곡선을 이어갈 수 있다. 도함수의 양의 하한과 \(C^2\) root family의 정칙성이 유지되고, 이 기울기 상계가 \(b=1\)에 도달하지 않음을 보장하기 때문이다. 마찬가지로 \(\lambda_0\cdot1.489886688502190\times10^{-4}<b_0\)이면 \(\lambda=0\)까지 이어갈 수 있다.

\textbf{증명.} (33)--(35)는 같은 \(\alpha\)에서 좌표 선분별로 (32)를 적분한 결과다. (36)은 implicit function theorem과 분모의 양의 하한에서 따른다. continuation은 위 기울기 상계와 compact 내부 구간의 정칙성으로 닫힌다. \(\square\)

네 모서리의 HH--birth 혼합 contrast

\[
\mathcal I=\Delta T(1,1)-\Delta T(1,0)-\Delta T(0,1)+\Delta T(0,0)
\]

에는 이번 \(\lambda\)-도함수 폭만으로

\[
\boxed{|\mathcal I|\le1.958485175177039\times10^{-9}\ {\rm K}}
\tag{37}
\]

가 성립한다. 다음 단위는 이 contrast와 \(\partial_{\lambda b}\Delta T\)를 원 mixed carry에서 직접 포함하여 상호작용의 방향과 크기를 정한다. 식 (35)--(37)의 exact rational 계산은 \nolinkurl{theory/PARAMETER_COROLLARIES.json}에 보존했다.

\hypertarget{uxad6cuxd604-uxc5f0uxacb0uxacfc-uxd655uxbcf4-uxbc94uxc704}{%
\section{구현 연결과 확보 범위}\label{uxad6cuxd604-uxc5f0uxacb0uxacfc-uxd655uxbcf4-uxbc94uxc704}}

\begin{enumerate}
\def\labelenumi{\arabic{enumi}.}
\tightlist
\item
  저장된 원 centre/box source Jet에서 실제 hull의 \(F,q\) 값·gradient와 whole-X signed Hessian을 회수한다.
\item
  PHYS09/14 \(U,V\)와 PHYS11/14 photon carry를 같은 \(\theta\)에 결박한다. 원 binary64 leaves는 exact rational로 해독한다.
\item
  식 (5)의 \(s_\theta\), (18)의 \(\Delta A_{1,\theta}\)를 먼저 구성하여 gas movement와 photon sector에 공유한다.
\item
  식 (9)--(11)로 \(k_{w\mathrm{He}}\)와 그 방향 도함수를 회수하여 식 (15),(16)의 source-time derivative를 닫는다.
\item
  식 (6),(19),(20)으로 \(r_\theta\)를 조립하고, 원 direct stage derivative 및 보존된 carry와의 포함 교차를 기록한다.
\item
  원 \(H_{gg},H_{g\theta}\)와 PHYS16 \(d\)로 식 (1)의 \(f_\theta\)를 포함하고 같은 \(C,B,S,q\)로 \(d_\theta\)를 얻는다.
\item
  같은 \(\Pi_F\)와 \(u,d,d_\theta\)로 식 (27) 또는 (28)을 적용한다. 실제 산술 출력에서 구간, 부호 margin, 남은 dependency width를 기록한다.
\item
  Wolfram/xAct는 식의 algebra·tensor contraction을 검산하고, 별도 검토는 원 source 구조·정규화·입력 연결과 실제 산술을 확인한다.
\end{enumerate}

이 파일의 모든 핵심 유도는 완결됐다. PHYS17의 실행 수치와 최종 채택 상태는 해당 원 result·review 파일에 결박한다. 새 strong no-go 주장은 제안하지 않는다.
